\section{Materials and Methods}

\subsection{Study Design and Datasets}

The development source was Coswara. The indexed release contained 2,746
participants; 2,114 participants with resolved positive or negative labels
contributed 19,024 recordings to supervised analyses. The remaining 632
participants (5,688 recordings) were retained in data audits but excluded from
model fitting and evaluation. Each participant may contribute several
prompted recordings, including heavy and shallow cough, deep and shallow
breathing, counting, and sustained vowels. All recordings from a participant
were assigned to one partition. Rows with unresolved labels remained in the
data-quality audit but were excluded from supervised evaluation.

The external target contained 8,331 COUGHVID recordings: 285 positive and
8,046 negative under the processed \texttt{status\_SSL} analytic label, for a
known-label prevalence of 3.42\%. Each external recording had a unique
identifier; COUGHVID does not provide reliable participant linkage for this
cohort, so the external analysis unit is a recording rather than a participant.
No
COUGHVID sample was used for source feature ranking, model fitting, fusion
selection, or source threshold selection. Because COUGHVID contains coughs,
external evaluation was restricted to the Coswara cough branch.

The two datasets do not implement identical outcome ascertainment. Coswara's
analytic labels originate from participant COVID-19 status fields, whereas the
COUGHVID target mapping uses the processed \texttt{status\_SSL} field after the
project's predefined quality and label rules. The raw COUGHVID \texttt{status}
field was available for fewer recordings and disagreed with \texttt{status\_SSL}
for 175 recordings with both fields resolved; a raw-status sensitivity analysis
remains outstanding. Label-definition mismatch is therefore a plausible
contributor to transfer error.

\subsection{Audio Quality and Preprocessing}

Audio paths, participant identifiers, modalities, labels, durations, and
quality flags were indexed before modeling. Loading converted supported source
formats to mono waveforms at 16~kHz and recorded decoding failures. Quality
checks identified 23,539 valid, 597 mostly silent, 390 corrupt, and 186 short
recordings in the indexed corpus. Conventional feature extraction removed
leading and trailing low-energy regions at a 35-dB threshold, localized an
active region by frame-level root-mean-square energy, peak-normalized the
signal, and center-cropped or padded cough, breathing, and speech to 4, 5, and
5~s, respectively. Time--frequency calculations used a 400-sample window,
160-sample hop, and 64 mel bands from 20 to 8,000~Hz. Each experiment declared
whether it used all decodable samples or only quality-passing samples; the rule
was fixed before feature computation and applied unchanged outside training.

\subsection{Acoustic Representations}

The main conventional representation concatenated three sources of
recording-level information: the existing acoustic bank, the ComParE 2016
functionals, and the INTERSPEECH 2010 Paralinguistic Challenge functionals.
The latter two banks were extracted with openSMILE
\cite{eyben2010opensmile,schuller2010is10,schuller2013compare}.
The resulting merged table contained 10,140 numeric candidates: 6,377
ComParE-derived columns, 1,586 IS10-derived columns, and 2,177 project acoustic
columns. LightGBM gain importance \cite{ke2017lightgbm} was calculated separately in each available
modality on the original source-training partition, normalized within modality,
and averaged across modalities. The primary reduced bank retained the 800
highest-ranked columns; the same names and order were frozen for validation,
test, chronological stress testing, and COUGHVID transfer. This design prevents
target-label feature selection but means the chronological analysis is not a
fully nested re-selection experiment.

Two learned representations tested whether the portability result was specific
to conventional functionals. WavLM Base Plus encoded short cough segments and
pooled segment predictions to the participant or recording level. The
spectrogram branch used a convolutional front end followed by a bidirectional
gated recurrent unit (CNN--BiGRU), so local time--frequency patterns were
learned before temporal aggregation. Both branches used source validation for
checkpoint and threshold decisions and were evaluated on COUGHVID without
target-label retraining.

\subsection{Classifiers and Imbalance Handling}

The conventional bank evaluated regularized radial-basis SVC, LightGBM,
XGBoost, and CatBoost classifiers
\cite{cortes1995svm,ke2017lightgbm,chen2016xgboost,prokhorenkova2018catboost}.
Imbalance handling was restricted to the
training partition. For model configurations denoted by an SMOTE fraction,
synthetic minority examples were generated only after the training subset and
feature space had been fixed; validation and test distributions were not
resampled \cite{chawla2002smote}. Hyperparameters and feature count were selected from training and
validation evidence, never from test or external labels.

\subsection{Multimodal Fusion}

Each modality model produced a probability $p_m(x)$ for participant $x$. Three
fusion rules were considered. Uniform fusion used
\begin{equation}
    p_{\mathrm{mean}}(x)=\frac{1}{|\mathcal{M}_x|}
    \sum_{m\in\mathcal{M}_x}p_m(x),
\end{equation}
where $\mathcal{M}_x$ is the set of available modalities. Validation-weighted
fusion transformed each modality's validation AUPRC as
$r_m=\max(\mathrm{AUPRC}_m-0.5,0.01)$ and normalized it,
\begin{equation}
    w_m=\frac{r_m}{\sum_j r_j},\qquad
    p_{\mathrm{weighted}}(x)=\sum_m w_m p_m(x).
\end{equation}
Stacked fusion fitted class-balanced L2 logistic regression only on aligned
validation predictions,
\begin{equation}
    p_{\mathrm{stack}}(x)
    =\sigma\!\left(\beta_0+\sum_m\beta_m p_m(x)\right),
\end{equation}
where $\sigma$ is the logistic function. The implementation consumes modality
probabilities directly; it does not first apply a logit transform. Fitted
coefficients and the validation-selected balanced-accuracy threshold were then
frozen for test evaluation. Consequently, stack coefficients were learned,
not hard-coded, and test labels did not determine them.

\subsection{Evaluation Protocols}

Four protocols were distinguished.

\begin{enumerate}
    \item \emph{Participant-disjoint internal evaluation} separated source
    participants across train, validation, and test partitions.
    \item \emph{Time-stratified participant evaluation} retained participant
    separation while distributing calendar periods across partitions. This
    tests a more calendar-balanced internal scenario.
    \item \emph{Early-to-late evaluation} trained on an earlier calendar block,
    selected models on source validation data, and evaluated later
    participants. The frozen development feature bank described above was used;
    feature selection was not refitted solely inside the early block. A reverse
    late-to-early analysis was exploratory.
    \item \emph{External transfer} trained and selected the cough branch on
    Coswara and applied it unchanged to COUGHVID.
\end{enumerate}

The protocol-specific selected models are not automatically paired estimands.
The four-number descriptive ladder is therefore supplemented by
model-and-modality-matched cough comparisons and paired tests only where the
same participants have predictions from both models.

Figure~\ref{fig:study-design} summarizes the development and evaluation
boundaries. In particular, the chronological and external analyses reuse the
source-developed representation and source-validation decisions; they do not
select a new model from the later or external labels.

\begin{figure}[!t]
\centering
\includegraphics[width=0.99\linewidth]{fig01_study_design.pdf}
\caption{Study design. Coswara supplied multimodal source development data;
COUGHVID supplied a cough-only external target. Participant separation,
training-only resampling and ranking, and source-validation model and threshold
selection were maintained throughout.}
\label{fig:study-design}
\end{figure}

\subsection{Confounding and Shift Analyses}

Metadata-only logistic models used three predefined input sets: symptoms,
demographic and recording-protocol variables, and their safe combined set.
Permutation importance grouped encoded columns by their source concept.
Full retraining with permuted labels tested whether the observed metadata
performance could arise from software leakage.

A source-versus-target classifier and its predicted-probability band measured
dataset separability and limited overlap. The reported fraction
outside the source probability band is a diagnostic, not a formal proof of a
causal positivity violation.

Temporal feature stability was measured by the Jaccard index between early and
late top-800 sets,
\begin{equation}
    J(F_E,F_L)=\frac{|F_E\cap F_L|}{|F_E\cup F_L|}.
\end{equation}

\subsection{Performance and Statistical Analysis}

Discrimination was summarized by AUROC and AUPRC. Threshold-dependent results
included sensitivity, specificity, precision, balanced accuracy, and F1. The
threshold was chosen on validation data by balanced accuracy unless a stated
fixed-sensitivity analysis was performed. Calibration was summarized by Brier
score, negative log likelihood, and expected calibration error (ECE):
\begin{align}
    \mathrm{Brier} &= \frac{1}{n}\sum_{i=1}^{n}(p_i-y_i)^2,\\
    \mathrm{ECE} &= \sum_{b=1}^{B}\frac{|I_b|}{n}
    \left|\operatorname{acc}(I_b)-\operatorname{conf}(I_b)\right|.
\end{align}
The reliability-bin formulation follows established probability-calibration
analysis \cite{niculescu2005calibration}.

For the target-only recalibration sensitivity analysis, COUGHVID was divided
before fitting into 2,082 calibration recordings and 6,249 disjoint evaluation
recordings. Platt scaling and isotonic regression were fitted only on the
calibration subset and assessed only on the evaluation subset. This analysis
tested whether probability rescaling could improve calibration and
threshold-dependent performance; it was not a new external test because target
labels were used to fit the recalibrator.

Bootstrap intervals resampled the declared analysis unit. For external AUROC
differences, source participants and target recordings were sampled
independently; repeated Coswara cough recordings were kept within participant
clusters. DeLong comparisons were used only for paired
predictions on the same labeled participants \cite{delong1988auc}. Decision-curve net benefit at
threshold probability $p_t$ was
\begin{equation}
    \mathrm{NB}(p_t)=\frac{TP}{n}-\frac{FP}{n}\frac{p_t}{1-p_t}.
\end{equation}
This definition compares the model with treat-all and treat-none strategies
across threshold probabilities \cite{vickers2006dca}.
The selected internal system, matched cough transfer, and representation-family
transfer formed the primary reliability analyses in this report. Subgroups,
context controls, reverse temporal testing, and mechanism analyses were
post-development exploratory analyses; their $p$-values were not interpreted
as a family of confirmatory tests.
